\documentclass[11pt,a4paper]{article}
\usepackage[utf8]{inputenc}
\usepackage[margin=1in]{geometry}
\usepackage{fancyhdr}
\usepackage{titlesec}

% Header setup to match the textbook layout
\pagestyle{fancy}
\fancyhf{}
\fancyhead[L]{82}
\fancyhead[R]{\textsc{Hornbill}}
\renewcommand{\headrulewidth}{0pt}

\setlength{\parindent}{0pt}
\setlength{\parskip}{1em}

\begin{document}

made at home by soaking the beans, grinding them and straining the water. (111 words)

A summary is usually one-third the length of the original passage. This is about half.

Now think of what we can omit to make the summary more brief as shown below.

\begin{quote}
The soybean leguminous plant which grows in all kinds of soil and climate yields beans, sprouts and a variety of processed food items and dairy alternatives and is also used to make candles and bio-diesel.

Rich in protein, vitamins, minerals and fibres, it has a low fat and cholesterol content. It lowers LDL levels and reduces the risk of coronary heart disease.

Soymilk which is lactose-free is available as flavoured milk and agrees with people allergic to ordinary milk. It can be made at home by soaking, grinding and straining soybean. (90 words)
\end{quote}

Try reducing it further to about 72 words.

\begin{quote}
Soybean, a legume, growing in a variety of soil and climatic conditions, yields beans, sprouts and a variety of food items and is used in making candles and bio-diesel.

Rich in protein, vitamins, minerals and fibres, it is low in cholesterol and fat. It lowers LDL levels and reduces the risk of coronary heart disease. Soymilk, lactose-free, is available flavoured and taken by people allergic to milk. It can also be made at home. (74 words)
\end{quote}

Notice that we have phrases in apposition: `a legume', between commas; present participles: `growing' to effect reduction. Instead of `it is rich in...' we have used `rich in...' and postponed the main verb in the sentence. Almost all the main points have been covered.

Read the text below and summarise it.

\begin{center}
    {\Large\bfseries Green Sahara} \\[0.3em]
    {\large\itshape The Great Desert Where Hippos Once Wallowed}
\end{center}

The Sahara sets a standard for dry land. It's the world's largest desert. Relative humidity can drop into the low single digits. There are places where it rains only about once a

\end{document}
\documentclass[11pt,a4paper]{article}
\usepackage[utf8]{inputenc}
\usepackage[margin=1in]{geometry}
\usepackage{fancyhdr}
\usepackage{textcomp}

% Header setup to match the textbook layout
\pagestyle{fancy}
\fancyhf{}
\fancyhead[L]{\textsc{Summarising}}
\fancyhead[R]{83}
\renewcommand{\headrulewidth}{0pt}

\setlength{\parindent}{2em}
\setlength{\parskip}{0.5em}

\begin{document}

\noindent century. There are people who reach the end of their lives without ever seeing water come from the sky.

Yet beneath the Sahara are vast aquifers of fresh water, enough liquid to fill a small sea. It is fossil water, a treasure laid down in prehistoric times, some of it possibly a million years old. Just 6,000 years ago, the Sahara was a much different place.

It was green. Prehistoric rock art in the Sahara shows something surprising: hippopotamuses, which need year-round water.

``We don't have much evidence of a tropical paradise out there, but we had something perfectly liveable,'' says Jennifer Smith, a geologist at Washington University in St. Louis.

The green Sahara was the product of the migration of the paleo-monsoon. In the same way that ice ages come and go, so too do monsoons migrate north and south. The dynamics of earth's motion are responsible. The tilt of the earth's axis varies in a regular cycle --- sometimes the planet is more tilted towards the sun, sometimes less so. The axis also wobbles like a spinning top. The date of the earth's perihelion --- its closest approach to the sun --- varies in a cycle as well.

At times when the Northern Hemisphere tilts sharply towards the sun and the planet makes its closest approach, the increased blast of sunlight during the north's summer months can cause the African monsoon (which currently occurs between the Equator and roughly 17\textdegree N latitude) to shift to the north as it did 10,000 years ago, inundating North Africa.

Around 5,000 years ago the monsoon shifted dramatically southward again. The prehistoric inhabitants of the Sahara discovered that their relatively green surroundings were undergoing something worse than a drought (and perhaps they migrated towards the Nile Valley, where Egyptian culture began to flourish at around the same time).

``We're learning, and only in recent years, that some climate changes in the past have been as rapid as anything underway today,'' says Robert Giegengack, a University of Pennsylvania geologist.

As the land dried out and vegetation decreased, the soil lost its ability to hold water when it did rain. Fewer clouds

\end{document}
\documentclass[11pt,a4paper]{article}
\usepackage[utf8]{inputenc}
\usepackage[margin=1in]{geometry}
\usepackage{fancyhdr}

% Header setup to match the textbook layout
\pagestyle{fancy}
\fancyhf{}
\fancyhead[L]{84}
\fancyhead[R]{\textsc{Hornbill}}
\renewcommand{\headrulewidth}{0pt}

\setlength{\parindent}{2em}
\setlength{\parskip}{0.5em}

\begin{document}

\noindent formed from evaporation. When it rained, the water washed away and evaporated quickly. There was a kind of runaway drying effect. By 4,000 years ago the Sahara had become what it is today.

No one knows how human-driven climate change may alter the Sahara in the future. It's something scientists can ponder while sipping bottled fossil water pumped from underground.

``It's the best water in Egypt,'' Giegengack said --- clean, refreshing mineral water. If you want to drink something good, try the ancient buried treasure of the Sahara.

\begin{flushright}
    \textsc{Joel Achenback} \\
    \textit{Staff Writer, Washington Post}
\end{flushright}

\end{document}

